% BHU PYQ -- Transcribed & Corrected
% Paper: PHYSE21: Energy Conversion and Storage for Sustainable Development (Mid Term)
% Exam: B.Sc. II Semester Mid Term Examination, 2024-25 / 2025 (NEP)
% Category: Skill Enhancement Course (SEC)
% Department: Physics

\documentclass[11pt,a4paper]{article}
\usepackage[a4paper,top=2.5cm,bottom=2.5cm,left=2.8cm,right=2.8cm,headheight=14pt]{geometry}
\usepackage{amsmath,amssymb}
\usepackage[T1]{fontenc}
\usepackage[utf8]{inputenc}
\usepackage{lmodern,microtype,fancyhdr,enumitem,hyperref,lastpage,parskip,tfrupee}

\hypersetup{
    colorlinks=true,
    urlcolor=black,
    linkcolor=black,
    pdftitle={PHYSE21: Energy Conversion & Storage (Mid Term) -- BHU Sem II 2024-25 (NEP SEC)}
}

\pagestyle{fancy}
\fancyhf{}
\renewcommand{\headrulewidth}{0.4pt}
\renewcommand{\footrulewidth}{0.4pt}
\fancyhead[L]{\small \textit{Banaras Hindu University}}
\fancyhead[R]{\small \textit{PHYSE21: Energy Conversion \& Storage (Mid Term)}}
\fancyfoot[C]{\small Page \thepage\ of \pageref{LastPage}}

\newcommand{\pts}[1]{\hfill \textit{[#1]}}

\begin{document}

\begin{center}
  {\large\bfseries BANARAS HINDU UNIVERSITY}\\[2pt]
  {\normalsize B.Sc. (Hons.) --- Semester II Mid Term Examination, 2024-25 (NEP)}\\[6pt]
  \rule{0.8\linewidth}{0.5pt}\\[6pt]
  {\large\bfseries PHYSICS (SEC)}\\[4pt]
  {\large\bfseries Paper Code: PHYSE21 : Energy Conversion and Storage for Sustainable Development}\\[4pt]
  {\normalsize Time: 1 Hour\hfill Maximum Marks: 20}
\end{center}

\rule{\linewidth}{0.4pt}

\smallskip

\noindent \textit{Instructions: Answer all questions.}

\bigskip

\noindent \textbf{Question 1.} Describe Biochemical energy conversion, and Solar energy conversion technologies.\pts{4}

\medskip

\noindent \textbf{Question 2.} Briefly describe the two major categories of thermal energy storage technologies.\pts{3}

\medskip

\noindent \textbf{Question 3.} In a water storage tank containing 1 metric ton of water, temperature of the water initially was 28~$^\circ$C and final temperature is 50~$^\circ$C. How much heat energy is stored?\pts{3}

\medskip

\noindent \textbf{Question 4.} Assume that you have five 100W lamps, ten 40W tube-lights, five 100W ceiling fans, one 400W refrigerator, and one 200W TV in your house.
\begin{enumerate}[label=(\alph*)]
  \item What is the minimum electricity load requirement for your house?
  \item Assume further that your electricity provider charges ₹5 per electricity unit. If all your appliances run at 50\% capacity about 8 hours a day on average, how much should be your monthly electricity bill?
  \item What is your annual electricity consumption in electricity units?
  \item Suppose that you can buy solar panel from the market in multiples of 1 kW panels and it costs ₹70,000 for each 1 kW panel. Assuming a 1 kW panel would produce 1400 electricity units on average in a year, how much money will you save per year if you install 2kW solar panels and in how many years you are going to recover your upfront installation cost?
  \item The PM Surya Ghar Yojana currently provides subsidy upto ₹78,000 (₹30,000 per kW up to 2 kW and ₹18,000 per kW for additional capacity up to 3 kW). Consider the case where you receive the eligible subsidy and calculate in how many years you will break even in this case.
  \item Would your savings increase if your house was in Jaipur instead of Varanasi?
\end{enumerate}\pts{6}

\medskip

\noindent \textbf{Question 5.} Under National Green Hydrogen Mission, how does India aim to source green Hydrogen? Name two domestic demands where green Hydrogen are expected to be utilized?\pts{2}

\medskip

\noindent \textbf{Question 6.}
\begin{enumerate}[label=(\alph*)]
  \item What are the broad objectives of the Standards \& Labelling Scheme of BEE?
  \item State True or False: \textit{Voltas's model-XYZ Air-conditioner received 4 star rating from BEE in 2024. In 2025, the same model received 5 star.}
\end{enumerate}\pts{3}

\end{document}
